\section{Introduction}
\label{sec:introduction}

Evaluating a candidate network against the true objective -- by training it, through weight sharing, or by
querying a benchmark -- dominates the cost of neural architecture search~\citep{white2023insights,liu2023survey}.
Architecture and training hyperparameters are commonly searched in sequence, yet Zela et al.\ show that they
interact, so the best architecture under one hyperparameter setting need not be the best under
another~\citep{zela2018towards}. Searching both jointly avoids that risk at the price of a larger space.

NSGA-Net established multi-objective evolutionary search over architectures~\citep{lu2019nsganet}, and
NSGANetV2 added a learned accuracy predictor that screens most candidates cheaply, reserving full evaluation for
a promising few~\citep{lu2020nsganetv2}. Both use NSGA-II, whose variation operators treat genotype coordinates
independently of whether they interact.

A separate family of evolutionary algorithms is built around exactly that interaction structure. GOMEA and the
Parameter-less Population Pyramid (P3) learn a linkage tree of dependent variables from the population and
protect those groups during crossover~\citep{thierens2011optimal,goldman2014parameterless}. On combinatorial
benchmarks their advantage over structure-blind search grows with problem size and
structure~\citep{goldman2015fast,qiao2024gomearuntime}, and surrogate-assisted variants have been
demonstrated~\citep{dushatskiy2019csgomea,dushatskiy2021novelsurrogateassistedevolutionaryalgorithm,przewozniczek2026lympus}.
Whether the combination transfers to NAS, and specifically to the joint architecture-and-hyperparameter setting,
is open.

\paragraph{This work.} P3Net replaces NSGA-II with P3 and replaces the absolute regressor with a
\emph{relative, linkage-aware} surrogate $\hat\delta_F$: instead of predicting a candidate's accuracy from its
full encoding, it predicts the change in accuracy relative to a parent, restricted to the one linkage subset that
was modified, in the spirit of eLyMPuS~\citep{przewozniczek2026lympus}. Our central question is whether, under a
fixed and small budget of full evaluations, this combination yields a better accuracy/cost trade-off than NSGA-II
with an absolute surrogate. We test it on JAHS-Bench-201~\citep{bansal2022jahsbench} (three datasets) and
NAS-HPO-Bench-II~\citep{hirose2021nashpobenchii}, against NSGANetV2, matched engine and surrogate ablations, and
the baselines established for this joint setting~\citep{guerreroviu2021bagofbaselines}.

\paragraph{The answer is negative.} P3Net does not separate from uniform random search in any of the sixteen
(search space, budget) cells. The rest of the paper is about why.

\paragraph{Contributions.}
\begin{itemize}
    \item A controlled comparison of a surrogate-assisted linkage-learning EA against NSGA-II-based and
    Bayesian-optimisation baselines for joint NAS and HPO, with matched engine and surrogate ablations,
    $R=30$ seeds, per-cell Holm--Bonferroni correction, and an absolute metric (IGD+ against the exact oracle
    front) on the one benchmark that admits it.
    \item A mechanistic diagnosis of the null result: P3's geometric population-pyramid growth forces a rising
    share of the budget into uniform-random bootstrap, and this property travels with the engine rather than
    with either surrogate.
    \item Two isolation experiments on architecture-only NAS-Bench-201 showing the null survives removing the
    hyperparameter dimension, and that a reconstruction of the closest prior surrogate-assisted
    linkage-learning design~\citep{bartnik2026evolutionary} reproduces it.
    \item An evaluation harness in which every number reported here is regenerated from persisted raw runs
    and checked against the text automatically.
\end{itemize}
